\todo[inline, author=Prof]{將發明的網路命名，提供兩個命名}
\todo[inline, author=Wu, color=blue!40]{ 1. self-attentive bidirectional GRU decoder \\   2. Attention-aided bidirectional decoder }

\section{Data acquisition}
\todo[inline, color=green!40]{ (換句話說Makin論文)包含訊號處理、擷取spike的部分，放入 Figure~\ref{fig:s1_and_m1_spike_train} }

% \begin{multicols}{2}
The first dataset, \textbf{Nonhuman Primate Reaching with Multichannel Sensorimotor Cortex Electrophysiology} is provided by \cite{makin2018} and made public available at \url{https://zenodo.org/record/583331#.Xv2qFygzaUk} with license under \href{https://creativecommons.org/licenses/by/4.0/legalcode}{CC BY 4.0}.
It containes the recordings from two Rhesus macaque monkeys, named Indy and Loco.
Each monkey contributed a several sessions of data.
One session represents a consecutive experiment recording within several minutes.
The lengths of the sessions are not consistant, and mostly the durations are in the range of 5 to 20 minutes.
Two 96-channel silicon microelectrode arrays (Figure~\ref{fig:utah_array}) were chronically implanted into their sensorimotor cortex.
Each array consisted of 96 electrodes, and the distance between two adjacent electrodes is $400 \mu m$. 
All electrodes covered the area of size $4mm \times 4mm$.
For each monkey two arrays were implanted, 
one was laid in primary motor cortex (M1) and the other was in primary somatosensory cortex (S1).
Some sessions of the dataset were constituted by the recording from both arrays, 
whereas most of the sessions consisted of the recording from the M1 array only.
The sampling rate of both arrays is 24.4kHz. 
More details regarding recording cortex area,
sorted units numbers, 
and duration of each session can be found in Table \ref{tab:dataset_indy} and \ref{tab:dataset_loco}.
% \end{multicols}

\begin{figure}[H]
    \begin{center}
        \begin{subfigure}{0.6\textwidth}
            \captionsetup{width=0.5\textwidth}
            \includegraphics[width=1\textwidth]{Utah_array.jpg}
            \caption{Utah array}
            \caption*{\small{source by \url{ https://commons.wikimedia.org/wiki/File:Utah_array_pat5215088.jpg }}}
            \medskip
        \end{subfigure}
        \begin{subfigure}{0.48\textwidth}
            \captionsetup{width=1\textwidth}
            \includegraphics[width=1\textwidth]{channel_array_1.png}
            \caption{ Indy M1, Loco M1 and Loco S1 }
        \end{subfigure}
        \begin{subfigure}{0.48\textwidth}
            \captionsetup{width=1\textwidth}
            \includegraphics[width=1\textwidth]{channel_array_2.png}
            \caption{ Indy S1 }
        \end{subfigure}
        \captionsetup{width=0.5\textwidth}
        \caption{ Electrode array and the apatial arrangement of 96 channels }
        \label{fig:utah_array}
    \end{center}
\end{figure}

% \begin{multicols}{2}
The raw recording from the electrode arrays were filtered with an IIR bandpass filter (fourth-order Butterworth; $F_{pass}$=500Hz to 5000Hz).
Spikes were detected and sorted online by \cite{makin2018} with a custom software written in C++.
In Figure \ref{fig:raw_to_spike}, 
the upper subplot is the raw broadband signal and the bottom subplot is the filtered high-frequency signal and the detected spikes.
% \end{multicols}

\begin{figure}[H]
    \begin{center}
        \includegraphics[width=1\textwidth]{raw_to_spike.png}
        \captionsetup{width=0.5\textwidth}
        \caption{Raw data waveform and the detected spikes}
        \fontsize{10}{12}
        \label{fig:raw_to_spike}
    \end{center}
\end{figure}

% \begin{multicols}{2}
To boost the accuracy of the decoders, 
BCI researches use past neural signals to predict the current kinematic 
The term \textit{tap size} represents the number of the discrete history inputs \citep{nakagome2020empirical}, 
In our case it is the number of past time bins used to decode each kinematics.
As illustrated in Figure~\ref{fig:time_bin_and_prediction}, 
the input dimension is the number of sorted units, and the model has the tap size of 4.
Since it performs a continuous decoding senario, 
the kinematic in each time bin expect for the very first 4 ones in the recording will be predicted.
The input time bins used to predict the kinematics at time $t$ and the time bins used to predict at time $t-1$ happen to overlap,
and since in this case the tap size is 4, the number of overlapping time bins are 3.
% \end{multicols}

\todo[inline, author=Prof.]{ Figure~\ref{fig:time_bin_and_prediction} 與tap size參考以下論文的Figure 3 Inhibition of Long-Term Variability in Decoding Forelimb Trajectory Using Evolutionary Neural Networks With Error-Correction Learning}
\begin{figure}[H]
    \begin{center}
        \includegraphics[width=0.7\textwidth]{time_bin_and_prediction.png}
        \captionsetup{width=0.5\textwidth}
        \caption{Tap sizes}
        \label{fig:time_bin_and_prediction}
    \end{center}
\end{figure}
\todo[inline, author=Wu, color=blue!40]{ 先說明 tap size ，time bin 的使用方式  }

% \begin{multicols}{2}
Due to the spike detecting threshold they determined in different sessions, 
the number of sorted units in each session varied from 2 to 4.
But all sessions it preserved an unsorted \textit{hash} unit in each channel.
These units contained all neuron firing events which surpassed the voltage threshold but they did not match any of the operator defined templates.
\cite{todorova2014sort} showed that there exists some information in the hash unit thus keeping them is benefitial in the movement decoding tasks. 
As showed in Figure \ref{fig:indy_20160407_02_M1_ch_60}, 
we visualized the M1 spike train in channel 60 from session \textbf{indy\_20160407\_02}.
The black rasters are the spikes classified as hash unit, 
and this session has two more sorted units, as they were marked in red and blue rasters respectively.
Figure \ref{fig:indy_20160407_02_S1_ch_60} showed the spike train from S1 electrode array in the same session and same recording interval.
% \end{multicols}

\begin{figure}[H]
    \centering
        \begin{minipage}[t]{0.49\textwidth}
            \centering
            \includegraphics[width=1\textwidth]{M1_Spike_Train_Channel_060.png}
            \caption{M1 Spike Train}
            \label{fig:indy_20160407_02_M1_ch_60}
        \end{minipage}
        \begin{minipage}[t]{0.49\textwidth}
            \centering
            \includegraphics[width=1\textwidth]{S1_Spike_Train_Channel_060.png}
            \caption{S1 Spike Train}
            \label{fig:indy_20160407_02_S1_ch_60}
        \end{minipage}
\end{figure}
\todo[inline, author=Wu, color=blue!40]{ 再說明有些 session 有 M1 S1 兩個電極的資訊 }

% \begin{multicols}{2}
Spike firing rate was calculated by counting the spike numbers within non-overlapping 64ms time window.
For sessions contained both M1 and S1 electrode array, 
we concatenated the two spike trains.
As from Figure~\ref{fig:spike_train_to_firing_rate} the spike trains turned into firing rate.
For other experiments, different bin width were used to obtain different firing rates.
% \end{multicols}

% \begin{figure}[H]
%     \centering
%         \begin{minipage}[t]{0.49\textwidth}
%             \centering
%             \includegraphics[width=1\textwidth]{s1_and_m1_spike_train.png}
%             \caption{M1 and S1 spike train}
%             \subcaption*{session indy\_20160411\_02}
%             \label{fig:s1_and_m1_spike_train}
%         \end{minipage}
%         \begin{minipage}[t]{0.49\textwidth}
%             \centering
%             \includegraphics[width=1\textwidth]{the_result_heatmap.png}
%             \caption{Firing rate}
%             \subcaption*{session indy\_20160411\_02}
%             \label{fig:the_result_heatmap}
%         \end{minipage}
% \end{figure}

\begin{figure}[H]
    \begin{center}
        \begin{subfigure}{0.49\textwidth}
            \includegraphics[width=1\textwidth]{s1_and_m1_spike_train.png}
            \caption{M1 and S1 spike train}
        \end{subfigure}
        \begin{subfigure}{0.49\textwidth}
            \includegraphics[width=1\textwidth]{the_result_heatmap.png}
            \caption{Firing rate}
        \end{subfigure}
        \caption{Spike train to firing rate}
        \subcaption*{session indy\_20160411\_02}
        \label{fig:spike_train_to_firing_rate}		
    \end{center}	
\end{figure}

\todo[inline, author=Wu, color=blue!40]{ 如何將 M1 與 S1 電極收到的 spike train 轉換成 firing rate，從 raster plot 變 heatmap。同時解釋 M1 與 S1 的 firing rate 是直接 concatenate 在一起。 }

% \begin{multicols}{2}
The heatmap of the firing rate within 50 time bins and the 6 kinematic variables are showed in Figure \ref{fig:Firing_rate_heatmap_and_all_kinematic_variables}.
For each type of the kinematic variable (position, velocity, acceleration) we will predict the trajectory among x and y-axis only without the z-axis.
The heatmap is showing an obvious relationship between the occurrance of movements and the intensity of spike firing.
Empty units in this session have been removed, as actual number of sorted units in each session is listed in Table \ref{tab:dataset_indy} and \ref{tab:dataset_loco}.
% \end{multicols}


\begin{figure}[H]
    \begin{center}
        \includegraphics[width=1\textwidth]{Firing_rate_heatmap_and_all_kinematic_variables.png}
        \captionsetup{width=0.5\textwidth}
        \caption{Firing rate and kinematic variables}
        \subcaption*{session indy\_20160927\_04}
        \label{fig:Firing_rate_heatmap_and_all_kinematic_variables}
    \end{center}
\end{figure}
\todo[inline, author=Wu, color=blue!40]{最後在將 firing rate heatmap 與 運動資訊疊圖比較 }


\section{Animal Pursuit task with point-to-point trajectory}
% \begin{multicols}{2}
\todo[inline, color=green!40]{(換句話說Makin論文)，包含運動軌跡、experimental setup}

The monkeys were trained to conduct reaching tasks with their arms in a transverse plane below their shoulder, 
and the targets were showing up one by one in a virtual reality environment.
Circular targets in 5mm virual radius were presented in Figure \ref{fig:panels} which the size is either $8 \times 8 $ square grid or $8 \times 17$ rectangular grid for monkey Indy,
and for monkey Loco square grid of size $6 \times 6$ was used.
% \end{multicols}

\begin{figure}[H]
    \begin{center}
        \centering
        \begin{subfigure}{0.8\textwidth}
            \centering\includegraphics[width=1\textwidth]{indy_20160411_02_panel.png}
            \label{fig:8_times_17}
            \caption{$8 \times 17$ rectangular grid} 
        \end{subfigure}

        \medskip
        \begin{subfigure}{0.4\textwidth}
            \centering\includegraphics[width=1\textwidth]{indy_20161025_04_panel.png}
            \label{fig:8_times_8}
            \caption{$8 \times 8 $ square grid} 
        \end{subfigure}
        \begin{subfigure}{0.4\textwidth}
            \centering\includegraphics[width=1\textwidth]{loco_20170214_02_panel.png}
            \label{fig:6_times_6}
            \caption{$6 \times 6$ square grid} 
        \end{subfigure}
        \caption[Three sizes of the panel]{Three sizes of the panel}
        \label{fig:panels}		
    \end{center}	
\end{figure}
% \todo[inline, author=Wu, color=blue!40]{ 3個panel的長寬尺寸還沒調成一致 }

% \begin{multicols}{2}
Only when the monkeys' fingertips reached the acceptance zone of size $7.5mm  \times 7.5mm$ which were centered on each target for 450ms would the system consider a valid touch and move on to another cue.
Figure \ref{fig:X-Y_trajectory_in_VR_panel} 
% \end{multicols}

\begin{figure}[H]
    \begin{center}
        \includegraphics[width=0.6\textwidth]{X-Y_trajectory_in_VR_panel.png}
        \captionsetup{width=0.5\textwidth}
        \caption{Fingertip Trajectory in 2D Virtual Panel}
        \label{fig:X-Y_trajectory_in_VR_panel}
    \end{center}
\end{figure}

% \begin{multicols}{2}
The recorded tragectories in the x and y-axis as well as the target cues for both directions are in Figure \ref{fig:x_and_y_trajectory_and_cue},
where the dash lines indicate the actual tragectory of the fingertip and the solid vertical gray lines indicate the time point when target cues altered.
% \end{multicols}

\begin{figure}[H]
    \begin{center}
        \includegraphics[width=1\textwidth]{x_and_y_trajectory_and_cue.png}
        \captionsetup{width=0.5\textwidth}
        \caption{Fingertip trajectory with respect to time}
        \label{fig:x_and_y_trajectory_and_cue}
    \end{center}
\end{figure}
\todo[inline, author=Wu, color=blue!40]{ 這邊真實紀錄到的軌跡都用虛線表示，是為了跟 \ref{sec:qualitative_analysis} 章 prediction、actual 軌跡對時間作圖保持一致， }

% \begin{multicols}{2}
They used a mirror to reflect visual input from a projector so that visual cues appeared in the plane of the reaching hand, 
so we must multiply a constant $-10$ to both x and y-axis recordings in order to transit the fingertip coordinate to the target coordinate with the unit in micrometer ($mm$).
An opaque barrier lay between the eyes of the monkeys and their fingertips,
so that the monkeys were not able to see their hands directly during experiments.
The experimental environment is showed in Figure \ref{fig:monkey_experiment}.
Fingertip position was recorded simultaneously by a 6-axis electromagnetic position sensors,
with a 250Hz sampling rate.
The fingertip velocity and acceleration were then computed afterwards by discrete derivative, 
which is simply the differences between two samples divided by the sampling period.
% \end{multicols}


\begin{figure}[H]
    \begin{center}
        \includegraphics[width=0.8\textwidth]{monkey_experiment.png}
        \caption{Experimental setup}
        \label{fig:monkey_experiment}
    \end{center}
\end{figure}


\begin{comment}
\section{Animal center-out task}

% \begin{multicols}{2}
The second dataset is provided by \cite{dyer2017cryptography} and available at \url{https://github.com/KordingLab/DAD}.
Only 2 sessions from monkey C was considered in \cite{makin2018}, both of them are center out movements as showed in Figure~\ref{fig:xy_plane_trajectory_chewie}.
% \end{multicols}

\begin{figure}[H]
    \begin{center}
        \begin{subfigure}{0.48\textwidth}
            \includegraphics[width=1\textwidth]{xy_plane_tragectory_session_Chewie_10032013.png}
        \end{subfigure}
        \begin{subfigure}{0.48\textwidth}
            \includegraphics[width=1\textwidth]{xy_plane_tragectory_session_Chewie_12192013.png}
        \end{subfigure}
        \caption{2D trajectories of monkey C}
        \label{fig:xy_plane_trajectory_chewie}
    \end{center}
\end{figure}

% \begin{multicols}{2}
\noindent 8 magenta squares are the targets for the monkey to reach, and black dots are the actual recordings.
The monkey held a handle to control a computer cursor in a 2D virtual plane.
Unlike grid reaching task experiments, this dataset is constituted by trials.
At the beginning of each trail the subject must move the cursor into the center area of size $2cm \times 2cm$.
After holding for 500 to 1500 ms, another 2cm target in the outer position will appear.
8 outer targets are at a 8cm radial distance from the center.
For monkey C, an auditory 'go' cue was provided right after the outer target showed up.

Apparently there exists a deviation between the targets coordinate and the cursors.
The center point of the cursor's coordinate in these two sessions are inconsistent as well.
and the transition matrix is not provided.
Fortunately the target points are not required when implementing hand movement decoders.
% \end{multicols}

\end{comment}

\section{Network architecture}
% \begin{multicols}{2}

We have combined the temporal self-attention module, layer normalization and bidirectional GRU together.
Figure~\ref{fig:Lin_Attention} showed our model structure.
Where $T$ is the tap size, $u$ is the GRU hyperparameter and is 256 in one-way, 512 in bidirectional GRU.
Since this model only performes self-attention along temporal dimension, 
the dimension of $u$ will be reduced to $d_a$, 
is chose to be 125 for simplicity, 
no matter it is one-way or bidirection GRU.
$r$ is the hops of attention we intend to extract from the same sequence, 
and it is set as 5 in decoding position and 2 in decoding velocity and acceleration.
% \end{multicols}
\begin{figure}[H]
    \begin{center}
        \includegraphics[height=0.5\textheight]{Lin_Attention.png}
        \captionsetup{width=0.5\textwidth}
        \caption{Model with temporal attention}
        \fontsize{10}{12}
        \label{fig:Lin_Attention}
    \end{center}
\end{figure}

\subsection{Recurrent Neural Networks}
\label{sec:RNN}
Figure~\ref{fig:unfold_rnn} illustrates how basic RNNs \textit{recur}.
The weight in each RNN cell remains the same during the process.
Only after it finishes a whole input sequence can the weight of RNN be adjusted by the output loss.
This diagram explained how the input and output mapped by the same RNN cell in each time step.

\begin{figure}[H]
    \begin{center}
        \includegraphics[width=1\textwidth]{unfold_rnn.png}
        \captionsetup{width=0.5\textwidth}
        \caption{Unfolded basic recurrent neural network}
        \caption*{\small{source by \url{https://en.wikipedia.org/wiki/File:Recurrent_neural_network_unfold.svg} licensed under \href{https://creativecommons.org/licenses/by-sa/4.0/deed.en}{CC BY-SA 4.0} }}
        \label{fig:unfold_rnn}
	\end{center}
\end{figure}

In our study, we used the model name Gated Recurrent Unit (GRU) as neural decoder.
The structure of a GRU cell is showed in Figure \ref{fig:GRU_cell}.

\begin{figure}[H]
    \begin{center}
        \includegraphics[width=0.7\textwidth]{GRU_cell.png}
        \captionsetup{width=0.5\textwidth}
        \caption[GRU Cell]{Mecahnism inside a GRU Cell}
        \caption*{\small{adapted from \url{ https://en.wikipedia.org/wiki/File:Gated_Recurrent_Unit,_type_2.svg } licensed under \href{https://creativecommons.org/licenses/by-sa/4.0/deed.en}{CC BY 4.0} }}
        \label{fig:GRU_cell}
    \end{center}
\end{figure}

% \begin{multicols}{2}

\begin{align}
    h_t &= \text{GRU} \left( x_t, \ h_{t-1}   \right)
\end{align}

\noindent and the four equations in the GRU cell:

\begin{align}
    % \label{eq:gru_cell}
    % \begin{split}
        r_t &=  \sigma (W_r h_{t-1} + U_r x_t + b_r) \label{eq:reset_gate} \\
        z_t &= \sigma (W_z h_{t-1} + U_z x_t + b_z) \label{eq:update_gate} \\
        \tilde{h}_t &= \phi(W (r_t \circ h_{t-1}) + U x_t + b) \label{eq:h_hat} \\
        h_t &= z_t \circ h_{t-1} + (1-z_t) \circ \tilde{h}_t \label{eq:h_out}
    % \end{split}
\end{align}

\noindent and the two activation functions:

\begin{align}
    \sigma (x) &= \frac{1}{1+e^{-x}} \label{eq:sigmoid} \\
    \phi (x) &= \frac{ e^{x}-e^{-x} }{ e^{x}+e^{-x} } \label{eq:tanh}
\end{align}

Where $r_t$ is the \textbf{Reset Gate} and $z_t$ is the \textbf{Update Gate}.
Notation $\circ$ means elementwise multiplication for two tensors of same shape.
Operator $\sigma$ is the sigmoid function (Equation \ref{eq:sigmoid}),
which squashes the input value from ($-\infty$, $\infty$) to (0,1).
$\phi$ is the hyperbolic tangent (tanh) function (Equation \ref{eq:tanh}),
squashing the input value from ($-\infty$, $\infty$) to (-1,1).
Two functions plotted in Figure \ref{fig:sigmoid_and_tanh}.
However both of them have drawbacks of saturation problem, 
e.g., the derivative of input value outside the range [-4, +4] will be nearly 0 after activations, 
making its gradient vanish and slowing down learninig process.
The solution about this issue will be addressed in Section~\ref{sec:layer_normalization}.
% \end{multicols}

\begin{figure}[H]
    \begin{center}
        \begin{subfigure}{0.48\textwidth}
            \includegraphics[width=1\textwidth]{sigmoid_and_tanh.png}
            \captionsetup{width=1\textwidth}
            \caption{Sigmoid and hyperbolic tangent}
        \end{subfigure}
        \begin{subfigure}{0.48\textwidth}
            \includegraphics[width=1\textwidth]{derivative_sigmoid_and_tanh.png}
            \captionsetup{width=1\textwidth}
            \caption{First derivative of sigmoid and hyperbolic tangent}
        \end{subfigure}
        \caption{Activation functions}
        \label{fig:sigmoid_and_tanh}
	\end{center}
\end{figure}

% \begin{multicols}{2}
At each times tep, a single GRU cell can take in two tensors, 
one is the current input tensor $x_t$ and another is the previous hidden state $h_{t-1}$, 
which is the output from last time step.
In the process of calculating each gate value (Equation \ref{eq:reset_gate}, 
\ref{eq:update_gate}) and $\tilde{h}_t$ (Equation \ref{eq:h_hat}),
hidden state from last time step ($h_{t-1}$) will be multiplied by unique weights ($W_r$, $W_z$, $W$),
and so does the input tensor ($x_t$) will be multiplied by its own weights ($U_r$, $U_z$, $U$).
After the multiplication, 3 biases ($b_r$, $b_z$, $b$) will be added to the results before entering the activation functions

A hyperparameter named \textbf{hidden dimension} or \textbf{number of hidden units} specified the dimensions for all hidden state in each time step, which is vector.
This hyperparameter also determines the number of each gates, 
i.e. with the setting of $u$ hidden units,
the shape of all hidden states should be $1 \times u$ and there will be $u$ reset gates, 
$u$ update gates in each cell.
Typically with the larger the hidden dimension, 
the more complex the model will be.
But highter complexity does not guarantee better performance since sophisticated models tend to overfit.

Figure \ref{fig:two_layer_GRU} demonstrated how hidden states are passed from one layer to another.
Where $T$ is the tap size of our input,
representing how many past time bins to be fed into the model just to decode the current kinematics.
The input data at each time step contains the firing rate of this moment $x_t$ and the past firing rate from $x_{t-n+1}$ to $x_{t-1}$.
Outputs of our model are the predicted kinematic variables for time bin $t$.
% \end{multicols}

\begin{figure}[H]
    \begin{center}
        \includegraphics[width=0.5\textwidth]{1-way_2-layer_GRU.png}
        \captionsetup{width=0.5\textwidth}
        \caption{Architecture of One-way Two-layer GRU}
        \fontsize{10}{12}
        \label{fig:two_layer_GRU}
    \end{center}
\end{figure}

% \begin{multicols}{2}
In the first GRU layer, $T$ bins of firing rate are fed chronologically.
Hidden state $h^{j}_{i}$ denotes the output tensor from the cell in timestep $i$ and layer $j$, and will be the input to the cell in timestep $i+1$ and layer $j+1$.
The size of $h^{j}_{i}$ should be $1 \times u$ since $u$ is a predefined parameter.
Ideally tensor $h^{2}_{t}$ should contain the information from the firing rate $x_{t-n+1}$ to $x_{t-1}$, thus by feeding it to some fully connect layers we can acquire the desired output.

The so called bidirectional RNN is made by simply adding another feeding of the same input sequence in reverse order.
As in Figure \ref{fig:bidir_two_layer_GRU}, but the size of the hidden state $h^{j}_{i}$ will be $1 \times 2u$,
since the first half of the this vector contains the information by feeding input sequence forward the last half by feeding backward.


\begin{align}
    \begin{split}
    \overrightarrow{h_t} &= \overrightarrow{\text{GRU}} \left( x_t  , \ \overrightarrow{h_{t-1}}  \right) \\
    \overleftarrow{h_t}  &= \overleftarrow{\text{GRU}}  \left( x_t  , \ \overleftarrow{h_{t+1}}   \right) \\
    h_t &= \overrightarrow{h_t} \oplus \overleftarrow{h_t}
    \end{split}
\end{align}

% \end{multicols}

\begin{figure}[H]
    \begin{center}
        \includegraphics[width=0.5\textwidth]{2-way_2-layer_GRU.png}
        \captionsetup{width=0.5\textwidth}
        \caption{Architecture of Bidirectional Two-layer GRU}
        \label{fig:bidir_two_layer_GRU}
    \end{center}
\end{figure}

It is the internal dynamics of the GRU to determine the degree of forget or keep the previous information in every time step.
In Section~\ref{sec:internal_dynamics_of_network} we provided figures to visualize these internal dynamics.


\subsection{Layer normalization}
\label{sec:layer_normalization}
% https://towardsdatascience.com/understanding-dataset-shift-f2a5a262a766
% \begin{multicols}{2}
    
A novel normalization method proposed by \cite{ba2016layer} called \textit{layer normalization} is designed to overcome the setbacks of batch normalization when it is implemented in RNNs.
Internal covariate shift is a problem caused by the variant distribution of the output of a previous layer,
which are used to be the input of its following layer.
It can be mitigated by fixing the mean and the variance of the summed inputs within each layer.
By the approach of normalization it can also provent the gradient vanishing problem from happening, which was brought up in Section~\ref{sec:RNN}.

Here is the computation of mean and variance from all the hidden units in the same layer:

\begin{align}
    \begin{split}
        \label{eq:normalization_term}
        \mu^l &= \frac{1}{H} \sum^{H}_{i=1} a^l_i \\
        \sigma^l &= \sqrt{ \frac{1}{H} \sum^H_{i=1} \left( a^l_i -\mu^l  \right)^2 }
    \end{split}
\end{align}

\noindent where $H$ denotes the number of hidden units in a layer, 
all the hidden units in a layer share the same normalization terms $\mu$ and $\sigma$.
Unlike batch normalization, 
layer normalization does not impose any constraint on the size of a mini-batch,
such that it cam be implemented to the model with batch size 1.

The original layer normalization methods is computed within a RNN cell.
In a standard RNN, 
the summed inputs in the recurrent layer are computed from the current input $x_t$ and previous vector of hidden states $h^{t-1}$, 
which are computed as $a^t = W_{hh} h^{t-1} + W_{xh} x^{t} $. 
Where $W_{hh}$ is the recurrent hidden to hidden weights and $W_{xh}$ are the input to hidden weights.
The layer normalizationed recurrent layer re-centers and re-scales its activations using the extra normalization terms similar to Equation \ref{eq:normalization_term}:

\begin{align}
    \begin{split}
        h^{t} &= f \left[ \frac{g}{\sigma^t} \odot \left( a^t- \mu^t \right) +b \right] \\
        \mu^t &= \frac{1}{H} \sum^{H}_{i=1} a^{t}_{i} \\
        \sigma^t &= \sqrt{ \frac{1}{H} \sum^{H}_{i=1} \left( a^t_i - \mu^t \right)^2 }
    \end{split}
\end{align}


\noindent where $\odot$ is the elementwise multiplication, 
and $b$, $g$ are defined as the bias and gain parameters of the same dimension as $h^t$.
$f$ is the activation function in the RNN cell, it could be a sigmoid or hyperbolic tangent function.

We define layer normalization as a function mapping $LN: \mathbb{R}^D \rightarrow \mathbb{R}^D$ with two set of adaptive parameters,
gains $\alpha$ and biases $\beta$:

\begin{align}
    \begin{split}
        LN(z; \ \alpha, \ \beta) &= \frac{(z-\mu)}{\sigma} \odot \alpha + \beta \\
        \mu &= \frac{1}{D} \sum^{D}_{i=1} z_i \\
        \sigma &= \sqrt{ \frac{1}{D} \sum^{D}_{i=1} \left( z_i- \mu \right)^{2} }
    \end{split}
\end{align}

When implementing it on GRU,

\begin{align}
    \begin{split}
        \begin{pmatrix}
            z_t \\
            r_t
        \end{pmatrix} &=  \sigma( W_h h_{t-1} ) + \sigma( W_x x_t ) \\
        \tilde{h}_t &= \text{tanh}( W x_t +  r_t \odot \left( U h_{t-1} \right) ) \\
        h_t &= \left( 1- z_t \right) h_{t-1} + z_t \tilde{h}_t
    \end{split}
\end{align}

\noindent layer normalization in GRU is applied as follows:

\begin{align}
    \begin{split}
        \begin{pmatrix}
            z_t \\
            r_t
        \end{pmatrix} &=   \sigma( LN( W_h h_{t-1} ; \ \alpha_1, \ \beta_1 ) ) + \sigma( LN(  W_x x_t ; \ \alpha_2 , \ \beta_2 ) ) \\
        \tilde{h}_t &= \text{tanh} (  LN( W x_t ; \ \alpha_3, \ \beta_3 ) +  r_t \odot LN( U h_{t-1} ; \ \alpha_4, \ \beta_4 ) ) \\
        h_t &= \left( 1- z_t \right) h_{t-1} + z_t \tilde{h}_t
    \end{split}
\end{align}

\noindent where $\alpha_i$ and $\beta_i$ are initialized as vector of zeros and ones, respectively.

The affect of the layer normalization is illustrated in Figure~\ref{fig:layer_norm_before_after},
where we simulated two extreme cases of hidden state distributions.
First one has a dense distribution but the mean is deviated from zero, 
making most of its values above $+4$, 
and they will be mapped approximately to zero after applied to in both sigmoid or hyperbolic tangent activation.
The second case has the mean of zero, but its broad distribution makes only a portion of its elements have their gradients non-zero after the activations.
% \end{multicols}

\begin{figure}[H]
    \begin{center}
        \begin{subfigure}{0.48\textwidth}
            \captionsetup{width=1\textwidth}
            \includegraphics[width=1\textwidth]{activation_histogram_before.png}
            \caption{Before LN}
        \end{subfigure}
        \begin{subfigure}{0.48\textwidth}
            \captionsetup{width=1\textwidth}
            \includegraphics[width=1\textwidth]{activation_histogram_after.png}
            \caption{After LN}
        \end{subfigure}
        \captionsetup{width=0.5\textwidth}
        \caption{LN comparison}
        \label{fig:layer_norm_before_after}
    \end{center}
\end{figure}

\subsection{Attention module}
\label{sec:attention_module}

\begin{figure}[H]
    \begin{center}
        \includegraphics[width=0.7\textwidth]{my_pipeline.png}
        \captionsetup{width=0.5\textwidth}
        \caption{ Attention-aided movement decoder }
        \label{fig:my_pipeline}
    \end{center}
\end{figure}

% \begin{multicols}{2}

The attention module is used to process the output features from its previous layers,
typically a RNN decoder in many cases \citep{zhang2020motor, hu2018novel}.
The structure of the whole decoding procedure is illustrated in Figure~\ref{fig:my_pipeline}.

\cite{lin2017structured} proposed a concise but effective attention module.
The output hidden states from the last layer of GRU will be the input of our attention module.
The hidden states matrix,
\begin{align}
    \begin{split}
        ( \bh_1, \ \bh_2, \cdots \ \bh_T) &= \bH \in \mathbb{R}^{T \times 2u }  \  \text{bidirectional RNN}  \\
           &= \bH \in \mathbb{R}^{T \times u} \ \ \text{one-way RNN}
    \end{split}
\end{align}
\noindent where $u$ is the number of hidden units and $ T $ is the length of each input sequence.
$T$ in our case it is the number of past time bins used to decode the current kinematic variables.
The matrix $\bH$ is the temporal features calculated by RNN in each time step, which is $T$ time steps in total.
Followed by the attention module, and it can give us the annotation vector, computed as
\begin{align}
    \label{eq:annotation_vector_calculation}
    \begin{split}
        \ba &= \text{softmax}( \bw_{s2} \ \text{tanh} ( \bW_{s1} \ \bH^{\intercal} ) )  \\
        \bm &= \ba \bH
    \end{split}
\end{align}
\noindent where weight matrix $\bW_{s1} \in \mathbb{R}^{d_a \times 2u} $ and vector $\bw_{s2}  \in \mathbb{R}^{ 1 \times d_a }$.
$d_a$ is an arbitrary hyperparameter.
The annotation vector $ \ba  \in \mathbb{R}^{ 1 \times T } $ will have its elements re-distributed by a softmax function $\mathbb{R}^K \rightarrow \mathbb{R}^K$,
which is defined as
\begin{align}
    \begin{split}
    \text{softmax}( \bz )_i &= \frac{ e^{ \bz_i } }{ \sum^{K}_{j=1} e^{ \bz_j } } \\ 
    \text{for} \ i &=1, \cdots , K \\ 
    \text{and} \ \bz&= \left( \bz_1, \cdots, \bz_K \right) \in \mathbb{R}^K
    \end{split}
\end{align}
\noindent such that the summation of the vector $ \ba $ is $1$.
The embedding vector $ \bm \in \mathbb{R}^{1 \times 2u} $ is the new features with enhanced temporal attribute.
Every single feature from each neural time bin has been annotated by an score.

% This explain fc in matrix form https://stackoverflow.com/questions/51608879/weights-matrix-final-fully-connected-layer

The vector representation may focues on only one specific part of the input sequence, 
and is expected to reflect only one \textit{aspect} in the sequence.
Multiple aspects can be extracted if we implement multiple hops of attention to the same input sequence.
Another hyperparameter, $r$ is designed to determine how many different parts can be extracted from the sequence.
The representations from Eqation \ref{eq:annotation_vector_calculation} is now changed to the matrix form in Eqation \ref{eq:annotation_matrix_calculation},
\begin{align}
    \label{eq:annotation_matrix_calculation}
    \begin{split}
        \bA &= \text{softmax}( \bW_{s2} \ \text{tanh}( \bW_{s1} \ \bH^{\intercal} ) )  \\
        \bM &= \bA \bH
    \end{split}
\end{align}
\noindent where the vector $ \bw_{s2}  \in \mathbb{R}^{ 1 \times d_a } $ turned into a matrix$ \bW_{s2} \in \mathbb{R}^{r \times d_a }  $,
and the embedding vector $ \bm \in \mathbb{R}^{1 \times 2u} $ became a matrix $\bM \in \mathbb{R}^{r \times 2u} $.

Eqation \ref{eq:annotation_matrix_calculation} is equivalent to a 2-layer MLP.
The first layer reduces the dimension from $2u$ to $d_a$, 
and the second layer reduces the dimension from $d_a$ to $r$, as illustrated in Figure~\ref{fig:multiply_attetion_vecotr_to_hidden_states}.

% \end{multicols}


\begin{figure}[H]
    \begin{center}
        \includegraphics[width=1\textwidth]{multiply_attetion_vecotr_to_hidden_states.png}
        \captionsetup{width=0.5\textwidth}
        \caption{Multiply attention vectors to hidden states}
        \fontsize{10}{12}
        \label{fig:multiply_attetion_vecotr_to_hidden_states}
    \end{center}
\end{figure}



\subsection{Penalization Term}
\label{sec:penalization_term}
% \begin{multicols}{2}
The purpose of developing such loss function in this attention module is to avoid repetitive information.
The embedding matrix $M$ should contains $r$ different aspects from the same input neural signal sequence.
But without further constraints, the attention module might always provide similar summation annotations for all the $r$ hops.
\cite{lin2017structured} proposed an approach to optimize the distribution of the annotations across the sequence,
named the \textit{penalization term} and it is defined as follows,

\begin{align}
    \label{eq:penalization_term}
    P &= {\norm{ \bA \bA^{\intercal} - \bI }_F}^2
\end{align}

\noindent where $A A^{\intercal} \in \mathbb{R}^{ r \times r} $, 
and $\norm{\bullet}_F$ denotes the Frobenius norm of a matrix, defined as
\begin{align}
    \label{eq:frobenius_norm}
    \norm{ \bX }_F &= \sqrt{ \sum^m_{i=1} \sum^n_{j=1} {|x_{ij}|}_2 }
\end{align}

% \end{multicols}

\begin{figure}[H]
    \begin{center}
        \includegraphics[width=1\textwidth]{penalization_term.png}
        \captionsetup{width=0.5\textwidth}
        \caption{Penalization term}
        \fontsize{10}{12}
        \label{fig:penalization_term}
    \end{center}
\end{figure}

% \begin{multicols}{2}

\noindent Figure~\ref{fig:penalization_term} showed that the summation of the squared value of every element in the matrix $\bA \bA^{\intercal} - \bI$ should be closed to $0$.
By implementing $P$ as a loss function,
its loss will be propagated backwards to alter the weights in the our attention module.

To further explain its concept,
let's start by considering two different annotation vectors $ \ba^i \in \mathbb{R}^{1 \times T} $ and $ \ba^j \in \mathbb{R}^{1 \times T}$.
Due to the fact that $ \text{softmax} $ is implemented,
the summation from all the element within the same vector is $1$.
The values in each vector can be considered as probability distribution.

In an extreme case, 
there should have the no overlapping between the probability distributions in these two vectors, 
thus the result from dot products $ \ba^i \cdot \ba^j $ should be $0$, and the calculation can be expressed as
\begin{align}
    0 < a_{ij} &= \sum^{n}_{k=1} a^i_k a^j_k < 1
\end{align}
\noindent where $a_{ij}$ is a non-diagonal element in matrix $\bA \bA^{\intercal}$, 
also should be equal to $0$.
$ a^i_k $ and $ a^j_k $ are the $k$-th elements in vector $ \ba^i$ and $ \ba^j $ severally.
The diagonal elements on the other hand,
will always be $ 1 $ since they are the dot products of two identity vectors,
$ \ba^i \cdot \ba^i $, i=1 $\sim$ r. 
That's why substracting an identity matrix from $A A^{\intercal}$ is needed when calculating $P$.
\cite{lin2017structured} claimed that in this way, 
each annotation vector $a^i$ is encouraged to focus on fewer parts of the sequence,
enabling each vector be focused on a single aspect. 

% \end{multicols}

\subsection{Visualization}
% \begin{multicols}{2}
Interpretation of the neural signals is one of the main topic in our study.
By taking advantage of attention module, we can visualize what the model \textit{sees.} \citep{xu2015show}
The term \textit{tap size} refer to the past time bins used to decode the current sample point.
In our main result, 20 tap size was used in decoding position and 10 tap size in decoding velocity and acceleration.
With the aid of annotation matrix from the attention module,
it is possible to depict the importance of the past time bins with annotations.
As showed in Figure~\ref{fig:visualization_vector}
% \end{multicols}

\begin{figure}[H]
    \begin{center}
        \includegraphics[height=0.2\textheight]{visualization_vector.png}
        \captionsetup{width=0.5\textwidth}
        \caption{Annotation visualization}
        \fontsize{10}{12}
        \label{fig:visualization_vector}
    \end{center}
\end{figure}

\section{Loss function and optimizer}
\label{sec:loss_function_and_optimizer}

% \begin{multicols}{2}
The first loss function for the predicted kinematic variables is mean square error (Equation \ref{eq:mean_square_error}),
which is widely used in regression problems.

\begin{align}
    \label{eq:mean_square_error}
    L_{MSE} &= \frac{1}{n} \sum^{n}_{i=1} \left( y_i - \hat{y}_i \right)^{2}
\end{align}

The second loss function is the penalization term explained in Section~\ref{sec:penalization_term}.
Summing them together with a coefficient for each loss function (Equation \ref{eq:loss_function_with_two_combined}),
$C_{1}$ and $ C_{2} $ are capable of weighing these two losses.

\begin{align}
    \label{eq:loss_function_with_two_combined}
    L &= C_{1} L_{MSE} + C_{2} P
\end{align}

Equation \ref{eq:adam_optimizer} is the Adam optimizer \citep{kingma2014adam} used in our research, 
a popular algothrim who combined AdaGrad \citep{duchi2011adaptive} and RMSProp \citep{tieleman2012lecture}.

\begin{align}
    \begin{split}
        \label{eq:adam_optimizer}
        m_t &= \beta_1 m_{t-1} + ( 1-\beta_1 ) \frac{\partial L_t}{\partial W_t} \\
        v_t &= \beta_2 v_{t-1} + (1- \beta_2) \left( \frac{\partial L_t}  {\partial W_t} \right)^2 \\
        \hat{m}_t &= \frac{m_t}{1- \beta^t_1 } \\
        \hat{v}_t &= \frac{v_t}{1-\beta^t_2} \\
        w_t &\leftarrow w_{t-1}- \eta \frac{ \hat{m}_t }{ \sqrt{\hat{v}_t} + \epsilon }
    \end{split}
\end{align}

\noindent Adam optimizer updates moving averages of the gradient $m_t$ and the squared gradient $v_t$,
where the hyperparameter $\beta_1, \ \beta_2$ control the exponential decay rates of these moving averages.

% \end{multicols}

\section{Performance evaluation}
% \begin{multicols}{2}
\todo[inline, color=green!40]{描述SNR、R2、single session}

To define the coefficient of determination $R^2$ and the signal-to-noise ratio (SNR), 
two vectors are required which are the observed value $ \by = \left[ y_1, \cdots , y_n  \right]^{\intercal} $ and predicted value $ \bf = \left[ f_1, \cdots , f_n \right]^{\intercal} $.
First we compute the mean of the observed data:
\begin{align}
    \bar{y} &= \frac{1}{n} \sum^{n}_{i=1} y_i
    \intertext{ then the total sum of squares, }
    SS_{tot} &= \sum^{n}_{i=1} \left( y_i - \bar{y} \right)^{2}
    \intertext{ the residual sum of squares, }
    SS_{res} &= \sum^{n}_{i=1} \left( y_i - f_i \right)^{2}
    \intertext{ with that the coefficient of determination will be }
    R^{2} &\equiv 1 - \frac{ SS_{res} }{ SS_{tot} } \label{eq:r_square}
    \intertext{ the SNR is defined by:}
    \text{SNR} &\equiv -10 \log_{10} (1 - R^{2})
\end{align}

\noindent and since the length of the testing data varied, weighted average is used to compute the mean performance across 47 sessions,
with the weights are the testing data length in each session.
% \end{multicols}